\documentclass{article}
\usepackage[utf8]{inputenc}
\usepackage{xspace}
\usepackage{fancyhdr}
\pagestyle{fancy}
\fancyhf{}
\rfoot{\thepage}
\lfoot{SISTCA}
\renewcommand{\headrulewidth}{0pt}

\title{\textbf{SISTEMAS COMPUTACIONAIS AVANÇADOS (SISTCA)
ADVANCED COMPUTING SYSTEMS}}

\author{Licenciatura em Engenharia Eletrotécnica e de Computadores}
\date{Instituto Superior de Engenharia do Porto 

Politécnico do Porto}


\begin{document}
\maketitle
\vspace{10cm}
\textbf{Lab Classes Script: }
\textbf{title}

\newpage

\textbf{Version control}
\vspace{1cm}


\begin{tabular}[h!]{|p{1.5cm}|p{1.5cm}|p{5cm}|p{4cm}|}
\hline
  \textbf{Version number}   & \textbf{Date issued} & \textbf{Authors} & \textbf{Update information} \\
  \hline
   \textbf{V1.0}  & month 2021 & Name surname (number) 
   
   name surname (number) 
   
   name surname (number) 
   
   name surname (number) 
   
   supervision: prof. name & Original version of the script \\
   \hline
   & & & \\
   \hline
   & & & \\
   \hline
   & & & \\
   \hline
   & & & \\
   \hline
\end{tabular}

\newpage

\tableofcontents
\newpage



\section{Introduction}
\subsection{Scientific/Technological Context}
(introduce the reader with the scientific and technological context, e.g. computer-aided 3D design/printing, computer vision/image processing, artificial intelligence/machine learning)
\subsection{Motivation for this topic/lab script}
(why the topic/technology is relevant in general terms and specifically in SISTCA/LEEC)
\subsection{Objectives (of this lab script}
(clearly identify the objectives of the lab script, i.e. summarizing what the student is supposed to learn (theoretical concepts, technology/platform), to do (hands-on, milestones), outlining the example/case-study and the challenge.
\subsection{Structure (of this lab script}
(summarize how the lab script is structured, in terms of sections: e.g. “The remainder of this document is organized as follows…”

\section{Theoretical (Scientific/technological background}
(this section will dig into the theoretical aspects underlying the technology addressed in the lab script; most bibliography should be referenced in this section)
\subsection{Main theoretical concepts/terminology}
(introduce the main theoretical concepts and the most relevant terminology, that will serve as a baseline for the rest of the document)
\subsection{State-of-the-art}
(short analysis of "state-of-the-art" technology related to the one addressd in this lab script, including a constructive comparison, identifying pros/cons; probably worth mentioning if the technology (software/hardware) is freely available, open-source, small/large community of users,…)

\section{Outlook of the technology}
(provide a snapshot of the main characteristics/functionality of the technology described in this lab script)
\subsection{Setup/installation (SW/HW)}
(describe the way to start using the technology, namely the setpup/installation SW/HW requirements, operating systems, development platforms,…)
\subsection{First steps}
(probably you can fit here something like the basic steps for a outputting a “hello world” message (broadly speaking))

\section{Tutorial/Functionality}
(this section is the core of the lab script, where all the “action” will happen; the student should be guided through a step-by-step execution, with supporting help/explanation of each functionality, towards a final objective, such as building a 3D model of a boat or producing a prototype hardware board)
\section{Exercises}
(this section may include one or more illustrative exercises, providing the corresponding resolution)
\subsection{Problem A}
\subsubsection{(re)solution A}
\subsection{Problem B}
\subsubsection{(re)solution B}
\section{Challenge}
(this section challenges students to do something different, to solve a specific problem, using the core technology/tool of this lab script; some hints may be provided, but not the (re)solution, which may be unveiled afterwards (e.g. one week later) or upon explicit solicitation from students that unsuccessfully attempted to solve the problem)
\section{References}
(here you should list all the bibliography that has been used as a base for writing this document and that have been effectively referenced in the text; the students are motivated to use referencing software such as the one embedded in MS Word, Mendeley, Zotero, Endnote, or any other of your choice; try to favour scientific-tecnological papers that have been published in high-qulity (or at least “trusted”) conferences and journals; both survey/position and research papers may be cited; avoid citing web sites, unless they hold web-apps or alike).
\section{Annexes}
(here you should include any technical documentation that you consider relevant for supporting the lab script but does not “fit into”, as you feel it would jeopardize the reading flow of the core text)

\end{document}
